\subsubsection*{Solution}
\paragraph*{Step-by-Step Derivation}

A necessary optical detail is not stated explicitly: two perpendicular standing waves form the usual separable square lattice only if the two beam pairs are mutually incoherent, normally by applying a small relative frequency shift. Otherwise, because all four beams have the same polarization, an additional interference term appears. This distinction is discussed in the optical-lattice literature; the standard four-beam square lattice uses independently detuned beam pairs \href{\detokenize{https://arxiv.org/abs/cond-mat/0403306}}{Blakie and Clark}.

The primary result below assumes this standard separable configuration.

\subparagraph*{1. Optical potential and lattice depth}

Let

\[
k=\frac{2\pi}{\lambda}, \qquad E_R=\frac{\hbar^2k^2}{2m} =\frac{2\pi^2\hbar^2}{m\lambda^2}
\]

be the laser wave number and single-photon recoil energy.

Take $E$ to be the peak electric-field amplitude of each traveling wave, and let $\alpha=\operatorname{Re}\alpha(\lambda)>0$. For a physical field written as

\[
\mathbf E_{\rm phys}(\mathbf r,t) =\operatorname{Re}\!\left[\boldsymbol{\mathcal E}(\mathbf r)e^{-i\omega t}\right],
\]

the cycle-averaged dipole potential is

\[
V_{\rm dip}(\mathbf r) =-\frac{\alpha}{4}\left|\boldsymbol{\mathcal E}(\mathbf r)\right|^2.
\]

This is the standard induced-dipole trapping potential \href{\detokenize{https://arxiv.org/abs/physics/9902072}}{Grimm, Weidemüller, and Ovchinnikov}.

The complex field of the pair propagating along $x$ is

\[
\mathcal E_x = 2E\,e^{-(y^2+z^2)/W^2}\cos(kx),
\]

and similarly

\[
\mathcal E_y = 2E\,e^{-(x^2+z^2)/W^2}\cos(ky).
\]

After averaging out the interference between the two independently detuned pairs,

\[
V_{\rm opt}(\mathbf r) = -V_0\left[ e^{-2(y^2+z^2)/W^2}\cos^2(kx) + e^{-2(x^2+z^2)/W^2}\cos^2(ky) \right],
\]

where the depth of each one-dimensional lattice is

\[
{V_0=\alpha E^2}.
\]

If a different electric-field convention is used, only this relation between $V_0$ and $\alpha E^2$ changes; the formulas below in terms of $V_0$ remain unchanged.

Define the dimensionless lattice depth

\[
s=\frac{V_0}{E_R}\gg1.
\]

The square-lattice spacing is

\[
a=\frac{\lambda}{2}=\frac{\pi}{k}.
\]

\subparagraph*{2. Tunneling energy}

Up to an irrelevant constant, the potential governing tunneling along either $x$ or $y$ is

\[
V_{\rm lat}(x)=V_0\sin^2(kx).
\]

The Hubbard tunneling matrix element is

\[
t=-\int d^3r\, w_i^*(\mathbf r) \left[ -\frac{\hbar^2\nabla^2}{2m}+V_{\rm opt}(\mathbf r) \right] w_j(\mathbf r),
\]

where $i,j$ are nearest-neighbor sites. In the Hamiltonian convention

\[
H_{\rm Hub} = -t\sum_{\langle ij\rangle,\sigma} \left(c_{i\sigma}^\dagger c_{j\sigma}+{\rm h.c.}\right) +U\sum_i n_{i\uparrow}n_{i\downarrow},
\]

$t$ is positive.

A local harmonic approximation by itself does not determine tunneling, because tunneling is controlled by the wavefunction under the barrier. In the stated deep-lattice limit, the controlled Mathieu/WKB asymptotic result is used. The under-barrier action is

\[
\frac{S}{\hbar} = \frac{1}{\hbar}\int_0^{\pi/k} \sqrt{2mV_0\sin^2(kx)}\,dx.
\]

Since $\sin(kx)\geq0$ on this interval,

\[
\frac{S}{\hbar} = \frac{\sqrt{2mV_0}}{\hbar} \int_0^{\pi/k}\sin(kx)\,dx = \frac{2\sqrt{2mV_0}}{\hbar k}.
\]

Using

\[
E_R=\frac{\hbar^2k^2}{2m},
\]

this becomes

\[
\frac{S}{\hbar}=2\sqrt{\frac{V_0}{E_R}}=2\sqrt{s}.
\]

Including the deep-well fluctuation prefactor gives

\[
{ t= \frac{4}{\sqrt{\pi}}E_R s^{3/4}e^{-2\sqrt{s}} }.
\]

This standard deep-lattice expression is also quoted in microscopic derivations of optical-lattice Hubbard models \href{\detokenize{https://journals.aps.org/prx/pdf/10.1103/PhysRevX.7.031057}}{APS, Phys. Rev. X 7, 031057}.

Corrections from the Gaussian envelope are of relative order $\lambda^2/W^2$ and are negligible under $W\gg\lambda$.

\subparagraph*{3. Harmonic Wannier functions}

Expand the optical potential near the central minimum $x=y=z=0$. Using

\[
\cos^2(kx)=1-k^2x^2+O(x^4),
\]

and

\[
e^{-2r^2/W^2} = 1-\frac{2r^2}{W^2} +O\!\left(\frac{r^4}{W^4}\right),
\]

one obtains

\[
V_{\rm opt} \simeq -2V_0 + V_0\left(k^2+\frac{2}{W^2}\right)(x^2+y^2) + \frac{4V_0}{W^2}z^2.
\]

Since $W\gg\lambda$, $2/W^2\ll k^2$, so

\[
V_{\rm opt} \simeq -2V_0 + V_0k^2(x^2+y^2) + \frac{4V_0}{W^2}z^2.
\]

Comparing with

\[
V_{\rm HO} = \frac12m\omega_\parallel^2(x^2+y^2) + \frac12m\omega_z^2z^2,
\]

gives

\[
\omega_\parallel = \sqrt{\frac{2V_0k^2}{m}} = \frac{2\sqrt{V_0E_R}}{\hbar},
\]

and

\[
\omega_z = \sqrt{\frac{8V_0}{mW^2}} = 2\sqrt{\frac{2V_0}{mW^2}}.
\]

The corresponding oscillator lengths are

\[
\ell_\parallel = \sqrt{\frac{\hbar}{m\omega_\parallel}} = \frac{1}{k}s^{-1/4},
\]

and

\[
\ell_z = \sqrt{\frac{\hbar}{m\omega_z}} = \sqrt{\frac{W}{2k}}\,s^{-1/4}.
\]

Thus the lowest-band Wannier orbital is approximated by

\[
w(\mathbf r) = \frac{ \exp\!\left[ -\frac{x^2+y^2}{2\ell_\parallel^2} -\frac{z^2}{2\ell_z^2} \right] }{ \pi^{3/4}\ell_\parallel\sqrt{\ell_z} }.
\]

\subparagraph*{4. On-site contact interaction}

For two distinguishable fermionic spin components, the low-energy contact coupling is

\[
g=\frac{4\pi\hbar^2a_s}{m}.
\]

The on-site Hubbard interaction is

\[
U=g\int d^3r\,|w(\mathbf r)|^4.
\]

For a normalized one-dimensional Gaussian of oscillator length $\ell$,

\[
\int_{-\infty}^{\infty}|w_\ell(q)|^4\,dq = \frac{1}{\sqrt{2\pi}\,\ell}.
\]

Therefore,

\[
\int d^3r\,|w(\mathbf r)|^4 = \frac{1}{ (2\pi)^{3/2}\ell_\parallel^2\ell_z }.
\]

Hence

\[
U = \frac{4\pi\hbar^2a_s}{m} \frac{1}{ (2\pi)^{3/2}\ell_\parallel^2\ell_z } = \sqrt{\frac{2}{\pi}} \frac{\hbar^2a_s}{ m\ell_\parallel^2\ell_z }.
\]

From the oscillator lengths,

\[
\ell_\parallel^2\ell_z = \frac{1}{k^2} \sqrt{\frac{W}{2k}}\, s^{-3/4},
\]

so

\[
\frac{1}{\ell_\parallel^2\ell_z} = \sqrt{2}\, \frac{k^{5/2}}{\sqrt W}\, s^{3/4}.
\]

Consequently,

\[
U = \frac{2}{\sqrt{\pi}} \frac{\hbar^2a_s}{m} \frac{k^{5/2}}{\sqrt W} s^{3/4}.
\]

Using $\hbar^2k^2/m=2E_R$,

\[
{ U= \frac{4}{\sqrt{\pi}}E_R \frac{ka_s}{\sqrt{kW}} s^{3/4} }.
\]

Since $k=2\pi/\lambda$, an equivalent form is

\[
{ U= 4\sqrt{2}\,E_R \frac{a_s}{\sqrt{\lambda W}} s^{3/4} }.
\]

The sign of $U$ is the sign of $a_s$: $a_s>0$ gives repulsion and $a_s<0$ attraction.

\subparagraph*{Coherent four-beam alternative}

If all four equally polarized beams are mutually coherent, their amplitudes must be added before squaring:

\[
V_{\rm opt} \simeq -\alpha E^2[\cos(kx)+\cos(ky)]^2.
\]

Defining the total well depth

\[
V_c=4\alpha E^2, \qquad s_c=\frac{V_c}{E_R},
\]

the lattice is rotated by $45^\circ$, its nearest-neighbor spacing is $\lambda/\sqrt2$, and the recoil energy along a nearest-neighbor direction is $E_R'=E_R/2$. The leading deep-lattice estimates then become

\[
t_c \simeq \frac{4}{\sqrt{\pi}}E_R' \left(\frac{V_c}{E_R'}\right)^{3/4} \exp\!\left[-2\sqrt{\frac{V_c}{E_R'}}\right],
\]

or

\[
t_c \simeq \frac{2}{\sqrt{\pi}}E_R (2s_c)^{3/4}e^{-2\sqrt{2s_c}},
\]

and

\[
U_c = 2^{7/4}E_R \frac{a_s}{\sqrt{\lambda W}} s_c^{3/4}.
\]

These differ from the standard square-lattice result because the prompt does not specify the mutual coherence of the two standing-wave pairs.
